% Author: Adam | Module: AI Job Interview Simulator | Project: SmartCareer AI
\section{Architecture}
\label{sec:archi}

\subsection{Vue d'ensemble}

La figure~\ref{fig:archi} donne la structure complète du module. Le client ne
parle jamais directement au moteur d'inférence : tout passe par le routeur
FastAPI, qui exige un jeton d'authentification, puis par le service
d'orchestration. Le service est le seul point de décision : c'est lui qui
compose le prompt, interroge l'index documentaire, valide la sortie, tente une
réparation et, si nécessaire, choisit un repli.

\begin{figure}[H]
  \centering
  \resizebox{\textwidth}{!}{%
  \begin{tikzpicture}[font=\sffamily]
    % ---- bande client ---------------------------------------------------
    \node[adamext]  (m1)   at (0,0)      {Éditeur de CV\\{\scriptsize Module 1}};
    \node[adamsvc]  (page) at (4.2,0)    {AdamInterviewPage\\{\scriptsize route /interview}};
    \node[adamsvc]  (queue) at (8.4,0)   {File d'attente RxJS\\{\scriptsize une analyse à la fois}};
    \node[adamsvc]  (api)  at (12.6,0)   {interview-adam.api.ts\\{\scriptsize HttpClient + Bearer}};
    \node[adambox]  (rel)  at (12.6,-2.2){releaseOnUnload()\\{\scriptsize fetch \ldots keepalive}};

    % ---- bande serveur --------------------------------------------------
    \node[adamext]  (app)  at (0,-4.4)   {app/api.py\\{\scriptsize montage du routeur}};
    \node[adamsvc]  (rout) at (4.2,-4.4) {router.py\\{\scriptsize /interview/*}};
    \node[adamsvc]  (serv) at (8.4,-4.4) {service.py\\{\scriptsize orchestration}};
    \node[adamsvc]  (llm)  at (12.6,-4.4){llm\_adam.py\\{\scriptsize client Ollama}};
    \node[adamext]  (oll)  at (16.8,-4.4){Ollama :11434\\{\scriptsize mistral 7B q4\_K\_M}};

    % ---- bande ressources ----------------------------------------------
    \node[adamsvc]  (prom) at (0,-7.4)   {prompts/*\\{\scriptsize 4 prompts}};
    \node[adamsvc]  (rag)  at (4.2,-7.4) {rag/retriever\\{\scriptsize FAISS + MiniLM}};
    \node[adamsvc]  (fb)   at (8.4,-7.4) {fallbacks.py\\{\scriptsize replis}};
    \node[adamsvc]  (db)   at (12.6,-7.4){db\_adam.py\\{\scriptsize sessions, qa}};
    \node[adamdb]   (sql)  at (16.8,-7.4){SQLite\\{\scriptsize data/cv\_platform.db}};
    \node[adamsvc]  (cvb)  at (4.2,-9.6) {cv\_bridge.py\\{\scriptsize lecture seule M1}};
    \node[adamsvc]  (priv) at (0,-9.6)   {privacy.py + security.py\\{\scriptsize anonymisation, gabarit}};

    % ---- flèches --------------------------------------------------------
    \draw[adamflow] (m1) -- (page);
    \draw[adamflow] (page) -- (queue);
    \draw[adamflow] (queue) -- (api);
    \draw[adamflow] (api) -- (rel);
    \draw[adamflow] (api.south) |- node[mlab,pos=.25,above]{HTTPS Bearer} (rout.north);
    \draw[adamflow] (app) -- (rout);
    \draw[adamflow] (rout) -- (serv);
    \draw[adamflow] (serv) -- (llm);
    \draw[adamflow] (llm) -- node[mlab,above]{POST /api/generate} (oll);
    \draw[adamback] (llm.south) |- (rag.north) node[mlab,pos=.75,below]{requête RAG};
    \draw[adamflow] (serv.south) -- node[mlab,right]{prompt} (prom.north);
    \draw[adamflow] (prom.west) .. controls +(-1,0) and +(0,1) .. (priv.east);
    \draw[adamflow] (rag) -- (cvb);
    \draw[adamflow] (serv.south) |- (fb.north);
    \draw[adamflow] (rout.south) |- (db.north);
    \draw[adamflow] (db) -- (sql);
    \draw[adamback] (rel.east) -| node[mlab,pos=.25,above]{pagehide / ngOnDestroy} (rout.east);

    % ---- légende --------------------------------------------------------
    \node[lam, anchor=north west] at (-1.6,-10.9)
      {\textbf{Légende} : boîte pleine = code du Module 3 \quad
       boîte pointillée = code du Module 1 ou service externe \quad
       flèche pointillée = appel secondaire ou asynchrone};
  \end{tikzpicture}}
  \caption{Architecture générale du simulateur d'entretien et points de contact
    avec le Module~1.}
  \label{fig:archi}
\end{figure}

Deux propriétés structurelles méritent d'être soulignées.

\paragraph{Un seul point de sortie vers le modèle.} Tout appel à Ollama passe
par \file{llm\_adam.py}. Cette concentration permet d'imposer une seule
politique : sémacphore de concurrence, bornes de timeout, retries avec reculée
exponentielle, messages d'erreur français et, surtout, un contrôle de
confidentialité émis \emph{avant} la moindre octet envoyé
(\const{assert\_anonymized(prompt)} dans \file{llm\_adam.py:143}).

\paragraph{Le module ne connaît le Module~1 que par une interface de lecture.}
\file{cv\_bridge.py} est le seul fichier qui importe \file{db.py}, et uniquement
\const{get\_generation}. Un test (\file{test\_cv\_bridge.py}) remplace les
fonctions d'écriture du Module~1 par une fonction qui lève, ce qui rend une
écriture accidentelle impossible sans le savoir.

\subsection{Structure du backend}

\begin{small}
\begin{tabularx}{\textwidth}{@{}>{\cellflush}p{4.3cm}
                             >{\cellflush}X@{}}
  \toprule
  \textbf{Fichier} & \textbf{Rôle}\\
  \midrule
  \file{config\_adam.py} (208 l.) & Toutes les constantes et leur variable
    d'environnement. Aucun autre fichier ne lit \file{.env}.\\
  \file{router.py} (176 l.) & 8 routes FastAPI sous le préfixe
    \route{/\brk{}interview}, toutes derrière \const{current\_user}. Traduit les
    erreurs internes en \code{HTTPException} à message français.\\
  \file{service.py} (511 l.) & Orchestration : contexte, RAG, prompt, appel,
    validation Pydantic, réparation, repli, persistance, synthèse.\\
  \file{llm\_adam.py} (318 l.) & Client HTTP d'Ollama. Sémacphore, retries,
    typage des pannes, analyse de la sortie JSON, \const{is\_available()},
    \const{release()}.\\
  \file{schemas\_adam.py} (425 l.) & Deux familles : les modèles Pydantic
    (contrat d'API et validation de sortie) et les schémas JSON aplatis envoyés
    dans \var{format}.\\
  \file{db\_adam.py} (242 l.) & Deux tables SQLite, schéma idempotent,
    connexions sérialisées par un verrou.\\
  \file{privacy.py} (186 l.) & Anonymisation du CV, nettoyage des textes
    libres, recherche de données personnelles résiduelles.\\
  \file{security.py} (60 l.) & Neutralisation des marqueurs de gabarit,
    troncature propre, encadrement des contenus par délimiteurs.\\
  \file{fallbacks.py} (191 l.) & Replis déterministes : banque de questions,
    analyse heuristique sans note, consigne STAR sans IA.\\
  \file{cv\_bridge.py} (99 l.) & Lecture seule des CV du Module~1, condensat
    anonymisé, écart offre/profil par vocabulaire de 96 termes, point de
    branchement pour les offres.\\
  \file{prompts/\brk} (5 fichiers) & \file{base\_adam.py} (assemblage Mistral,
    budget, délimiteurs) et les trois prompts métier plus le prompt de
    réparation.\\
  \file{rag/\brk} (5 fichiers) & Découpage Markdown, plongeur CPU, index FAISS,
    retriever, CLI d'ingestion et de recherche.\\
  \file{knowledge/\brk} (6 fichiers) & Base documentaire en anglais, 1\,150 lignes.\\
  \file{tests/\brk} (9 fichiers) & 2\,368 lignes, 201 tests.\\
  \bottomrule
\end{tabularx}
\end{small}

Le paquet est autonome : \file{interview\_adam/\brk\_\_\_init\_\_\_.py} déclare
qu'il ne modifie aucun fichier du Module~1 et ne réutilise en lecture que
\file{config.py}, \file{db.py}, \file{auth.py} et \file{cv\_schema.py}.

\subsection{Structure du frontend}

Le frontend occupe \file{frontend/\brk{}src/\brk{}app/\brk{}interview-\brk{}adam/\brk} (16 fichiers,
2\,751 lignes) et se limite à quatre composants autonomes, chargés par route
lazy :

\begin{itemize}[leftmargin=1.6em,itemsep=1pt,topsep=2pt]
  \item \file{adam-\brk{}interview-\brk{}page.ts} (79 l.) : coquille de page. Un signal
    \const{mode} vaut \const{'setup'}, \const{'run'} ou \const{'history'}.
    C'est lui qui branche l'écouteur \file{pagehide} et qui appelle la
    libération de VRAM.
  \item \file{adam-\brk{}interview-\brk{}setup.ts} (110 l.) : formulaire de lancement,
    valeurs par défaut (6 questions, mixte, junior, français, sans CV) et
    sélecteur de CV alimenté par \route{GET /history}.
  \item \file{adam-\brk{}interview-\brk{}run.ts} (222 l.) : l'entretien lui-même, une carte
    par question, la file d'arrière-plan et la synthèse.
  \item \file{adam-\brk{}interview-\brk{}history.ts} (84 l.) : historique, relecture en
    lecture seule, suppression avec confirmation.
\end{itemize}

Deux fichiers sans composant complètent l'ensemble : \file{models.ts} (175 l.),
qui transpose en TypeScript les modèles Pydantic du backend, et
\file{labels.ts} (40 l.), qui centralise les libellés français et le formatage
des notes.

\paragraph{État par question.} Chaque carte est décrite par deux cellules
indépendantes \const{EvalCell} et \const{ImproveCell}, dont l'état appartient à
\const{\{'idle' | 'queued' | 'running' | 'done' | 'error'\}}. Les valeurs par
défaut sont des constantes gelées partagées, pour ne pas recréer un objet vide
par question à chaque rendu.

\paragraph{La file d'attente.} Le composant de déroulement maintient trois
champs privés : une liste \const{queue} de travaux, un champ \const{running}
jouant le rôle de jeton d'exclusion mutuelle, et un drapeau \const{destroyed}.
Le texte de la réponse est \emph{figé à l'entrée en file} : attendre ne peut
pas faire analyser un texte que le candidat a modifié entre-temps.
La méthode \const{drain()} refuse d'entrer si le composant est détruit, si un
travail tourne déjà, ou si la file est vide ; sinon elle dépile un travail, le
marque \const{running} et s'abonne à l'appel HTTP.

Le point de conception le plus important est l'utilisation de
\const{finalize()} et non de \const{complete()} pour libérer le jeton
(\file{adam-\brk{}interview-\brk{}run.ts:175} et \file{:189}). Un flux RxJS en erreur
n'émet \emph{jamais} \const{complete} : avec \const{complete}, le premier échec
laissait \const{running} occupé indéfiniment et toutes les analyses restantes
resteront en attente. Ce bogue, trouvé en navigateur réel, est détaillé en
section~\ref{sec:bugs}.

\paragraph{Libération de la VRAM côté client.} Deux chemins coexistent, parce
que les deux situations sont différentes. Pour une navigation interne Angular
(la page est encore vivante), \const{ngOnDestroy()} appelle
\code{api.release()} via \const{HttpClient}. Pour un rechargement ou une
fermeture d'onglet, \const{HttpClient} est annulé par le navigateur pendant la
phase de déchargement du document : on utilise alors un \code{fetch} natif avec
\code{\{keepalive: true\}} (\file{interview-\brk{}adam.api.ts:79}), et comme
l'intercepteur Angular ne s'applique pas à \code{fetch}, l'en-tête
\code{Authorization} est posé à la main. Un drapeau \const{unloadFired} garantit
un appel unique.

\subsection{Modèle de données}
\label{sec:modeles}

Le module écrit dans le même fichier SQLite que la plateforme
(\file{data/\brk{}cv\_platform.db}), mais avec sa propre connexion et son propre
verrou : \file{db.py} n'est ni importé ni modifié. Le schéma est idempotent
(\code{CREATE TABLE IF NOT EXISTS}) et appliqué à la première utilisation ; le
mode WAL n'est pas activé, parce qu'il changerait le fichier partagé avec le
Module~1. Un \code{timeout=30} sur la connexion évite les erreurs
«~database is locked~» quand les deux modules écrivent le même jour.

\begin{small}
\begin{tabularx}{\textwidth}{@{}>{\cellflush}p{3.9cm}
                             >{\cellflush}p{2.2cm}
                             >{\cellflush}X@{}}
  \toprule
  \textbf{Colonne} & \textbf{Type} & \textbf{Remarque}\\
  \midrule
  \multicolumn{3}{@{}l}{\textbf{\route{interview\_sessions}}}\\
  \code{id} & TEXT PK & hexadécimal de 32 caractères (\code{uuid4().hex})\\
  \code{user\_id} & TEXT & \code{REFERENCES users(id) ON DELETE CASCADE}\\
  \code{job\_title}, \code{company}, \code{job\_description} & TEXT &
    contexte saisi, déjà nettoyé\\
  \code{seniority}, \code{question\_type}, \code{language} & TEXT & valeurs
    contrôlées par Pydantic à l'entrée\\
  \code{n\_questions} & INTEGER & nombre réellement retenu\\
  \code{model} & TEXT & nom du modèle ; \code{NULL} si la session est un repli\\
  \code{status} & TEXT & \code{in\_progress} ou \code{finished}\\
  \code{cv\_generation\_id} & TEXT & référence au CV du Module~1, sans copie\\
  \code{summary\_json} & TEXT & synthèse figée à la clôture\\
  \code{created\_at}, \code{finished\_at} & TEXT & horodatages ISO~8601 UTC\\
  \midrule
  \multicolumn{3}{@{}l}{\textbf{\route{interview\_qa}}}\\
  \code{id} & TEXT PK & identifiant technique\\
  \code{session\_id} & TEXT & \code{REFERENCES interview\_sessions(id)
    ON DELETE CASCADE}\\
  \code{idx} & INTEGER & numéro de question, à partir de 1\\
  \code{question}, \code{category}, \code{focus} & TEXT & question générée\\
  \code{expected\_json} & TEXT & barème attendu, liste de 2 à 4 points\\
  \code{answer} & TEXT & réponse du candidat, \code{NULL} tant que non analysée\\
  \code{score} & REAL & note globale ; \code{NULL} = non notée\\
  \code{criteria\_json}, \code{strengths\_json}, \code{weaknesses\_json} & TEXT
    & sortie d'évaluation\\
  \code{verdict} & TEXT & verdict rédigé\\
  \code{improved\_json} & TEXT & réponse améliorée + STAR + conseils\\
  \code{created\_at}, \code{answered\_at} & TEXT & horodatages\\
  \bottomrule
\end{tabularx}
\end{small}

Deux index accompagnent le schéma : \code{idx\_interview\_sessions\_user} sur
\code{(user\_id, created\_at DESC)} pour l'historique, et
\code{idx\_interview\_qa\_session} sur \code{(session\_id, idx)} pour la
relecture d'une session.

\paragraph{Relations et cloisonnement.} La seule clé étrangère sortante est
\code{interview\_sessions.user\_id $\rightarrow$ users(id)}, table du Module~1 ;
la suppression en cascade est donc double : supprimer un compte emporte les
sessions, et supprimer une session emporte ses réponses
(\code{interview\_qa.session\_id}). Mis à part cette clé, aucune table du
Module~1 n'est référencée : \code{cv\_generation\_id} est stocké comme un simple
texte, ce qui laisse le Module~1 libre de sa propre politique de suppression.

\paragraph{Isolation par utilisateur.} \emph{Toute} requête de lecture ou
d'écriture est bornée par \code{user\_id} dans sa clause \code{WHERE}
(\file{db\_adam.py:146}, \file{:217}, \file{:231}). Un identifiant de session
appartenant à un autre compte répond 404, jamais 403 : un 403 confirmerait
qu'une session existe, or l'identifiant est la seule chose qui protège le
contenu d'un entretien.

\subsection{Diagrammes de séquence}
\label{sec:seq}

Les quatre scénarios ci-dessous résument les échanges réels, tels qu'ils ont
été observés dans les journaux pendant les runs Playwright de
section~\ref{sec:tests}.

\subsubsection{Génération des questions}

\begin{figure}[H]
  \centering
  \resizebox{\textwidth}{!}{%
  \begin{tikzpicture}[font=\sffamily]
    \node[ll] (c) at (0,0)     {Candidat};
    \node[ll] (f) at (4.0,0)   {Angular};
    \node[ll] (s) at (8.0,0)   {service.py};
    \node[ll] (r) at (11.6,0)  {retriever};
    \node[ll] (l) at (15.2,0)  {Ollama};
    \node[ll] (d) at (18.4,0)  {SQLite};
    \foreach \x in {c,f,s,r,l,d} \draw[lifeline] (\x.south) -- ++(0,-6.6);

    \draw[msg] (c) -- node[mlab,above]{formulaire} (f);
    \draw[msg] (f) -- node[mlab,above]{POST /interview/questions} (s);
    \draw[msg] (s) -- node[mlab,above]{CV anonymisé, écarts} (d);
    \draw[msg] (s) -- node[mlab,above]{retrieve(`questions')} (r);
    \draw[rmsg] (r) -- node[mlab,above]{top-2 $\geq$ 0,20} (l);
    \draw[rmsg] (r) -- node[mlab,below]{contexte + sources} (s);
    \draw[msg] (s) -- node[mlab,above]{prompt + QUESTION\_SET\_SCHEMA} (l);
    \draw[rmsg] (l) -- node[mlab,below]{JSON validé par Pydantic} (s);
    \draw[msg] (s) -- node[mlab,above]{create\_session, attach\_questions} (d);
    \draw[rmsg] (s) -- node[mlab,below]{QuestionsOut} (f);
    \draw[rmsg] (f) -- node[mlab,below]{cartes + puces RAG} (c);
  \end{tikzpicture}}
  \caption{Génération des questions. Le trait pointillé rouge marque une
    réponse ; le RAG est interrogé \emph{avant} le modèle, et son absence ne
    change rien au reste du flux.}
  \label{fig:seq-questions}
\end{figure}

\subsubsection{Évaluation en arrière-plan, avec file}

\begin{figure}[H]
  \centering
  \resizebox{\textwidth}{!}{%
  \begin{tikzpicture}[font=\sffamily]
    \node[ll] (c) at (0,0)   {Candidat};
    \node[ll] (q) at (4.2,0) {File RxJS};
    \node[ll] (a) at (8.4,0) {API};
    \node[ll] (s) at (11.8,0){service};
    \node[ll] (l) at (15.4,0){Ollama};
    \foreach \x in {c,q,a,s,l} \draw[lifeline] (\x.south) -- ++(0,-7.2);

    \draw[msg] (c) -- node[mlab,above]{« Analyser Q1 »} (q);
    \draw[msg] (q) -- node[mlab,above]{POST /evaluate (Q1)} (a);
    \draw[msg] (a) -- (s);
    \draw[msg] (s) -- (l);
    \draw[msg] (c) -- node[mlab,above]{tape sa réponse à Q2} (c |- 0,-2.4);
    \draw[msg] (c) -- node[mlab,above]{« Analyser Q2 »} (q);
    \node[mlab, anchor=west] at (q.east |- 0,-3.6) {Q2 en file : \const{queued}};
    \draw[rmsg] (l) -- node[mlab,below]{Q1 évaluée (Q1 \const{done})} (s);
    \draw[rmsg] (s) -- (a);
    \draw[rmsg] (a) -- node[mlab,below]{feedback Q1} (q);
    \node[mlab, anchor=west] at (q.east |- 0,-5.4) {finalize $\rightarrow$ jeton
      libre $\rightarrow$ \const{drain()}};
    \draw[msg] (q) -- node[mlab,above]{POST /evaluate (Q2)} (a);
    \draw[msg] (a) -- (s);
    \draw[msg] (s) -- (l);
    \draw[rmsg] (l) -- (s);
    \draw[rmsg] (s) -- (a);
    \draw[rmsg] (a) -- node[mlab,below]{feedback Q2} (c);
  \end{tikzpicture}}
  \caption{Deux analyses successives avec \const{MAX\_CONCURRENT\_LLM\_CALLS=1}.
    Pendant que Q1 est en cours, la zone de saisie de Q2 reste utilisable ; Q2
    passe par \const{queued} puis \const{running} puis \const{done}.}
  \label{fig:seq-eval}
\end{figure}

\subsubsection{Amélioration à la demande}

\begin{figure}[H]
  \centering
  \resizebox{\textwidth}{!}{%
  \begin{tikzpicture}[font=\sffamily]
    \node[ll] (c) at (0,0)   {Candidat};
    \node[ll] (f) at (4.2,0) {Angular};
    \node[ll] (s) at (8.4,0) {service};
    \node[ll] (db) at (11.9,0){SQLite};
    \node[ll] (l) at (15.4,0){Ollama};
    \foreach \x in {c,f,s,db,l} \draw[lifeline] (\x.south) -- ++(0,-5.2);

    \draw[msg] (c) -- node[mlab,above]{clic « meilleure réponse »} (f);
    \draw[msg] (f) -- node[mlab,above]{POST /improve \{session, index\}} (s);
    \node[mlab, anchor=west] at (s.east |- 0,-1.5)
      {relit la question, la réponse et l'évaluation déjà en base};
    \draw[msg] (s) -- (db);
    \draw[rmsg] (db) -- node[mlab,below]{Q + réponse + faiblesses} (s);
    \draw[msg] (s) -- node[mlab,above]{prompt IMPROVEMENT\_SCHEMA} (l);
    \draw[rmsg] (l) -- node[mlab,below]{improved\_answer, star, tips} (s);
    \draw[msg] (s) -- node[mlab,above]{save\_improvement} (db);
    \draw[rmsg] (s) -- node[mlab,below]{ImprovementOut} (f);
    \draw[rmsg] (f) -- node[mlab,below]{bloc « Réponse améliorée »} (c);
  \end{tikzpicture}}
  \caption{L'amélioration relit l'évaluation déjà stockée : elle n'est pas
    appelée avec le prompt d'évaluation en mémoire, ce qui garantit qu'un
    rechargement de page ne change pas le résultat.}
  \label{fig:seq-improve}
\end{figure}

\subsubsection{Libération de la VRAM}

\begin{figure}[H]
  \centering
  \resizebox{\textwidth}{!}{%
  \begin{tikzpicture}[font=\sffamily]
    \node[ll] (n) at (0,0)    {Navigation};
    \node[ll] (f) at (4.4,0)  {Angular};
    \node[ll] (a) at (8.6,0)  {router.py};
    \node[ll] (l) at (12.4,0) {Ollama};
    \node[ll] (m) at (16.2,0) {Module 1};
    \foreach \x in {n,f,a,l,m} \draw[lifeline] (\x.south) -- ++(0,-4.6);

    \draw[msg] (n) -- node[mlab,above]{pagehide / ngOnDestroy} (f);
    \draw[msg] (f) -- node[mlab,above]{POST /interview/release (Bearer)} (a);
    \draw[msg] (a) -- node[mlab,above]{keep\_alive = 0} (l);
    \draw[rmsg] (l) -- node[mlab,below]{modèle déchargé} (a);
    \draw[rmsg] (a) -- node[mlab,below]{\{released, model\}} (f);
    \node[mlab, anchor=west] at (m.west |- 0,-4.0)
      {\route{GET /api/ps} : 0 Mio occupé};
    \draw[adamflow] (l) -- (m);
  \end{tikzpicture}}
  \caption{Libération explicite. Sans cet appel, \const{keep\_alive=5m} laisse
    Mistral resident jusqu'à 5 minutes et le Module~1 ne peut pas charger
    Qwen3-0.6B + LoRA.}
  \label{fig:seq-release}
\end{figure}

\subsection{Branchement avec le Module~1}
\label{sec:branchement}

Le contrat passé avec l'équipe était de ne toucher au code du coéquipier
qu'aux lignes strictement indispensables. Le résultat vérifié est le suivant.

\begin{small}
\begin{tabularx}{\textwidth}{@{}>{\cellflush}p{4.6cm}
                             >{\cellflush}p{2.0cm}
                             >{\cellflush}X@{}}
  \toprule
  \textbf{Fichier (propriété Module 1)} & \textbf{Lignes} & \textbf{Nature}\\
  \midrule
  \file{app/\brk{}api.py} & 56--57 & import du routeur, puis
    \code{app.include\_router(interview\_router)}. \emph{Rien d'autre} : une
    recherche de «~interview~» dans les 349 lignes du fichier ne renvoie que ces
    deux lignes.\\
  \file{frontend/\brk{}src/\brk{}app/\brk{}app.routes.ts} & 31--36 & une entrée de route
    \code{path: 'interview'} avec le garde \code{authGuard} existant et un
    \code{loadComponent} paresseux.\\
  \file{frontend/\brk{}src/\brk{}app/\brk{}pages/\brk{}workspace/\brk{}workspace.html} & 46--50 & un
    lien de menu «~Simuler un entretien~», qui réutilise la classe
    \code{.menu-item} du design system.\\
  \file{.env.example} & 9--60 & bloc de variables \file{ADAM\_*} et
    \file{MISTRAL\_*}. Les quatre variables du Module~1
    (\file{GEMINI\_API\_KEY}, \file{GEMINI\_MODEL},
    \file{GEMINI\_FALLBACK\_MODELS}, \file{LLM\_PROVIDER}) conservent leur valeur
    d'origine.\\
  \bottomrule
\end{tabularx}
\end{small}

\paragraph{Fichiers non touchés, vérifié.} Une recherche
\code{grep -i "interview|adam|mistral"} ne renvoie aucun résultat dans
\file{db.py}, \file{auth.py}, \file{generator.py}, \file{llm.py},
\file{prompts.py}, \file{cv\_schema.py}, \file{guards.py}, \file{agent.py},
\file{config.py} et \file{app/\brk{}main.py}.

\paragraph{Aucun CSS global ajouté.} Les quatre feuilles de style du module
consomment 20 jetons \code{var(-\brk{}-\brk{}\ldots)} déjà définis par le design system du
coéquipier dans \file{frontend/\brk{}src/\brk{}styles.scss}, et ne déclarent aucune variable
personnalisée. Le tableau des dépendances globales d'\file{angular.json}
contient toujours une seule entrée.

\paragraph{Ce qui reste à faire côté équipe.} Le \file{README.md} du dépôt
(79 lignes) ne mentionne pas encore le Module~3 : la section d'usage doit être
rédigée par l'équipe avec le reste de la documentation, ce qui correspond à
l'item~7 du plan de reprise.
